% ─────────────────────────────────────────────────────────────────────────────
% Section 1: Introduction
% Paper: "Silent Decisions Are Type Errors: Enforcing AI Accountability
%         via Linear Resource Types"
% Venue: IEEE Computer — Special Issue on AI Governance
% Deadline: 13 April 2026
% Authors: Daniel Lau Pereira Soares
% ─────────────────────────────────────────────────────────────────────────────

\section{Introduction}
\label{sec:intro}

Deployed AI systems now issue consequential decisions in milliseconds:
a credit application is rejected; a social-benefit disbursement is blocked.
Regulatory frameworks acknowledge this reality.
Brazil's Lei Geral de Proteção de Dados (LGPD), Article~20,
grants data subjects the right to request review of decisions based solely
on automated processing that affect their interests; on request, controllers
must provide clear and adequate information about the criteria and procedures,
subject to protection of trade and industrial secrets~\cite{brazil2018lgpd}.
The European Union AI Act requires human-oversight capabilities for high-risk
systems (Article~14) and, in defined cases, a right to explanation
(Article~86)~\cite{eu2024aiact}.
Neither instrument, by itself, prevents software from emitting a consequential
decision before the required evidence or review support exists.
A system may emit a denial, log it asynchronously, and silently drop the log.
No law is formally broken in that instant; the violation surfaces only later,
if at all.
We call this phenomenon a \emph{silent decision}:
an AI-generated act with legal consequences for an individual
that carries no integrity-authenticated evidence record.
Silent decisions are not a failure of intent---they are a failure of type.

The conventional response is instrumentation.
Policy engines such as Open Policy Agent~\cite{opa2024} and
Kyverno govern cloud infrastructure---pods, secrets, network policies---issuing
allow/deny verdicts over \emph{resources}, not over \emph{persons}.
Evidence of a decision's basis is a separable artifact, appended
asynchronously, omissible under load, and falsifiable after the fact.
Runtime assertions offer a second line of defense:
\texttt{assert(evidence != null)} at every callsite.
The analogy to memory safety is instructive.
In~C, \texttt{free()} can be paired with discipline and checked by analyzers---
yet dangling-pointer vulnerabilities accumulated for four decades because
discipline fails at scale.
Rust's ownership model solved the same problem by making the
\emph{absence of a deallocation owner} a compile-time error, not a runtime check.
The lesson generalizes:
no existing framework makes the absence of evidence a compile-time error
rather than a runtime omission.

Linear Logic, introduced by Girard in 1987~\cite{girard1987linear} and
connected to programming language semantics by Wadler~\cite{wadler1991linear},
treats propositions as \emph{resources}: they must be used exactly once,
can be neither duplicated nor discarded.
The connective $\multimap$ (linear implication) denotes a resource transformation
that \emph{consumes} its premise, and $\otimes$ (tensor product) denotes the
\emph{joint} consumption of two resources---not their independent conjunction.
We exploit this semantics to encode the accountability invariant directly
in the type system.
We use the type law $(E \otimes C) \multimap V$:
an \texttt{EvidenceToken}~($E$) and a \texttt{ComplianceToken}~($C$)
are jointly consumed to produce a Verdict~($V$).
Rust's ownership model is affine rather than strictly linear; Section~3
states the resulting limits explicitly.  In our implementation,
\texttt{EvidenceToken} is a \texttt{\#[must\_use]}
struct with no \texttt{Clone} or \texttt{Copy} derivation;
\texttt{Verdict::new()} takes both tokens by value, consuming them atomically.
Because \texttt{Blake3Hash}---the only evidence artifact---has no public
constructor, it cannot be forged: the sole path to a \texttt{Blake3Hash}
is through \texttt{EvidenceToken::consume()}.
A~silent decision is therefore not an illegal act that slips past a guard;
it is a \emph{type error}:
the compiler rejects it before a single byte of evidence is ever dropped.

This paper makes three contributions.
\begin{enumerate}
  \item \textbf{Linear Resource Types for AI Accountability (C1).}
    Building on linear resource discipline established for protocols
    and resources (Girard~\cite{girard1987linear}, Vault, typestate,
    session types; Section~\ref{sec:linear_related}), we instantiate it
    over AI accountability obligations---LGPD Art.~20 and EU AI Act
    Arts.~14 and~86: the core invariant
    $(E \otimes C_{\text{signed}}) \multimap V$ is a compile-time type
    constraint in Rust, verified on the companion
    repository~\cite{soares2026a, soares2026b}.

  \item \textbf{Decision-Evidence Causal Binding (C2).}
    We prove that, under the \emph{BTV}\footnote{BTV (\textit{Bound-To-Verdict}):
    a reference kernel implementing atomic accountability via linear types.
    Evidence and compliance tokens are structurally \emph{bound to} the
    verdict---the name is the type law $(E \otimes C) \multimap V$ in English.}
    type invariant, every materialized
    \texttt{Verdict} consumed \emph{exactly one} causally-prior
    \texttt{EvidenceToken}: a BLAKE3 digest of the decision context
    computed before the verdict existed.
    Linear ownership prevents duplication; the private constructor
    prevents forgery; both guarantees are static.

  \item \textbf{The Constitutional Enclosure Theorem (C3).}
    We state and prove that external Safe Rust callers cannot construct a
    BTV \texttt{Verdict} through the public API without both required tokens.
    The proof analyzes the crate's public construction paths and is
    machine-checkable: \texttt{cargo test} on the companion repository
    verifies all seventeen clauses in seconds.\
\end{enumerate}

\noindent
Section~2 situates the work against Policy-as-Code, audit logging, and
linear-type approaches.
Sections~3--5 develop the type system, state and prove the Constitutional
Enclosure Theorem, and report the implementation and benchmarks.
Section~6 discusses limitations and design trade-offs.
